%%%%%%%%%%%%%%%%%%%%%%%%%%%%%%%%%%%%%%%%%%%%%%%%%%%%%%%%%%%%%%%%%%%%%%%%%%%%
%% CS 550 Massive Data Mining and Learning -- Course Project Report
%% Template: ACM conference proceedings (acmart, sigconf), as suggested by
%% the project description. Compiled with pdflatex.
%%%%%%%%%%%%%%%%%%%%%%%%%%%%%%%%%%%%%%%%%%%%%%%%%%%%%%%%%%%%%%%%%%%%%%%%%%%%
\documentclass[sigconf,nonacm]{acmart}

%% Silence the ACM submission/rights banner -- this is a course project, not
%% an ACM submission.
\settopmatter{printacmref=false}
\setcopyright{none}
\renewcommand\footnotetextcopyrightpermission[1]{}
\pagestyle{plain}

% amsmath is already loaded by acmart. We do NOT load amssymb here because
% acmart's newtxmath already defines the symbols we need (incl. \Bbbk,
% \mathbb, etc.) and loading amssymb on top triggers a `\Bbbk already
% defined' error.
\usepackage{booktabs}
\usepackage{multirow}
\usepackage{enumitem}

%% ACM metadata (required fields, filled with course-local values)
\acmConference[CS 550]{CS 550 Massive Data Mining and Learning}{Spring 2026}{Rutgers University}
\acmYear{2026}
\copyrightyear{2026}

\begin{document}

\title{Seven Recommenders on MovieLens-1M: \\ Rating Prediction, Top-10 Ranking, and a Fairness Audit}
\subtitle{CS 550 Course Project, Option 1: Recommender System}

\author{Daize Dong}
\email{dd1376}
\affiliation{\institution{Rutgers University}\city{New Brunswick}\country{USA}}

\author{Jiawei Wu}
\email{jw2070}
\affiliation{\institution{Rutgers University}\city{New Brunswick}\country{USA}}

\begin{abstract}
We implement and compare seven recommender algorithms on the MovieLens-1M
dataset, covering the full range from trivial baselines to graph neural
networks: GlobalMean, BiasOnly, ItemKNN, UserKNN, biased Matrix Factorization
(MF), Neural Matrix Factorization (NeuMF), Bayesian Personalized Ranking
Matrix Factorization (BPR-MF), and LightGCN. All seven share a single
per-user 80/20 split and a single evaluation protocol, so differences come
only from modelling choices. We report the two rating-prediction metrics
required by the course specification (MAE, RMSE) and the four Top-10
ranking metrics (Precision@10, Recall@10, F$_1$@10, NDCG@10). As an
optional trustworthiness task, we additionally audit every model for
exposure fairness via Coverage, Gini, and long-tail share. Every model is
written from NumPy and PyTorch primitives; no recommender-system library
(Surprise, RecBole, LightFM, implicit, etc.) is imported. We find that the
training loss matters far more than the architecture on this dataset
(BPR-MF beats MF by $2.55\times$ on P@10 with the same scoring function),
that LightGCN does not beat BPR-MF at matched gradient budget, and that the
two best ranking models are also the two worst on fairness.
\end{abstract}

\maketitle

%-----------------------------------------------------------------------------
\section{Overview}
This project implements seven recommender models on the MovieLens-1M
dataset and compares them on two tasks required by the course
specification: (i)~rating prediction, evaluated with MAE and RMSE,
and (ii)~Top-10 item recommendation, evaluated with Precision@10,
Recall@10, F$_1$@10, and NDCG@10. The seven models cover four
families: simple baselines (GlobalMean, BiasOnly), memory-based
collaborative filtering (UserKNN, ItemKNN), matrix factorization
(MF, BPR-MF), a neural method (NeuMF), and a graph-based method
(LightGCN). Every model is implemented from scratch using NumPy and
PyTorch primitives; no recommender-system library is imported. We
use PyTorch layers and optimizers (Embedding, Linear, Adam) as
low-level building blocks only, per the course policy that forbids
using packaged recommender algorithms directly.

%-----------------------------------------------------------------------------
\section{Dataset and Preprocessing}
We use the MovieLens-1M dataset, the ``small'' movie recommendation dataset
listed in the project description. It contains $1{,}000{,}209$ integer
ratings on a 1 to 5 scale from $6{,}040$ users on $3{,}706$ movies, with each
user rating at least 20 movies (the 20-core condition).

\paragraph{Preprocessing.} We drop items with fewer than 5 ratings to
stabilize similarity estimation, leaving $N_u{=}6040$ users and
$N_i{=}3416$ items. User and item IDs are remapped to contiguous integer
ranges to allow dense embedding tables.

\paragraph{Split.} Following the course spec, for each user we uniformly
sample 20\% of their ratings as test and the remaining 80\% as train, with
random seed 42. Because every user has $\ge 20$ ratings in the raw data and
we do not drop users, every user has at least one training and at least one
test rating, so no special casing is needed. The final split is
$|\mathcal{D}_{\text{train}}|{=}799{,}721$ and
$|\mathcal{D}_{\text{test}}|{=}199{,}890$, with 95.16\% sparsity. Review
text is not used; timestamps are not used.

%-----------------------------------------------------------------------------
\section{Methods}
Notation: user $u$, item $i$, observed rating $r_{ui}$, global mean $\mu$,
user / item biases $b_u, b_i$, user and item latent factors
$p_u, q_i \in \mathbb{R}^d$.

\subsection{Baselines}
\textbf{GlobalMean} predicts $\hat{r}_{ui} = \mu$ for all pairs.
\textbf{BiasOnly} predicts $\hat{r}_{ui} = \mu + b_u + b_i$, where the
biases are learned by 15 alternating closed-form ridge updates with
$\lambda = 5$.

\subsection{Neighbourhood CF}
We use shrinkage-adjusted cosine similarity on mean-centred ratings. For
the item side:
\begin{equation}
s_{ij} = \frac{\sum_{u \in \mathcal{U}_{ij}} \tilde{r}_{ui} \tilde{r}_{uj}}{\|\tilde{r}_{\cdot i}\| \|\tilde{r}_{\cdot j}\|} \cdot \frac{|\mathcal{U}_{ij}|}{|\mathcal{U}_{ij}| + \alpha}, \qquad \tilde{r}_{ui} = r_{ui} - \bar{r}_u,\ \alpha = 100.
\end{equation}
We keep the top $k = 50$ neighbours per item. The prediction is:
\begin{equation}
\hat{r}_{ui} = \bar{r}_u + \frac{\sum_{j \in N_k(i) \cap \mathcal{I}_u} s_{ij} \tilde{r}_{uj}}{\sum_{j \in N_k(i) \cap \mathcal{I}_u} |s_{ij}|}.
\end{equation}
\textsc{UserKNN} works the same way but on the user side.

\subsection{Biased Matrix Factorization (MF)}
The prediction is $\hat{r}_{ui} = \mu + b_u + b_i + p_u^\top q_i$, trained
with L2-regularized MSE loss. We use $d = 64$, 20 epochs of Adam with
learning rate $10^{-2}$, batch size 8192, and $\lambda = 10^{-4}$.

\subsection{Neural Collaborative Filtering (NeuMF)}
NeuMF has two parallel branches: a GMF branch that computes
$p_u^g \odot q_i^g$, and an MLP branch on the concatenation
$[p_u^m \| q_i^m]$ with hidden sizes $(64, 32, 16)$, each layer followed by
ReLU and Dropout(0.1). The outputs of the two branches are concatenated
and passed through a linear layer. The original NeuMF uses binary
cross-entropy for implicit feedback; we use MSE instead because the task
here is explicit rating prediction.

\subsection{BPR-MF}
BPR-MF trains on triples $(u, i, j)$ where $i$ is an observed item and
$j$ is an unobserved item sampled uniformly with rejection against the
positive set. The loss is:
\begin{equation}
\mathcal{L}_{\text{BPR}} = -\sum_{(u,i,j)} \log \sigma(\hat{x}_{ui} - \hat{x}_{uj}) + \lambda \|\Theta\|^2, \qquad \hat{x}_{ui} = p_u^\top q_i + b_i.
\end{equation}
We use $d = 64$, 20 epochs, batch size 4096, Adam with learning rate
$10^{-2}$, and $\lambda = 10^{-5}$.

\subsection{LightGCN}
LightGCN removes feature transforms and non-linearities from the standard
GCN, keeping only weighted-average propagation on the user-item bipartite
graph. Let $\mathbf{A}$ be the adjacency matrix and
$\hat{\mathbf{A}} = \mathbf{D}^{-1/2} \mathbf{A} \mathbf{D}^{-1/2}$ be its
symmetric-normalized form. Each layer computes
$\mathbf{E}^{(\ell)} = \hat{\mathbf{A}} \mathbf{E}^{(\ell-1)}$, and the
final embedding averages over all layers:
\begin{equation}
\mathbf{E}^{*} = \frac{1}{L+1} \sum_{\ell=0}^{L} \mathbf{E}^{(\ell)}.
\end{equation}
Training uses the same BPR loss. We set $L = 3$ and $d = 64$. Because the
first LightGCN run used a much larger batch size than BPR-MF and therefore
many fewer gradient steps, we run LightGCN twice: once at 15 epochs with
batch size 16{,}384, and once at 20 epochs with batch size 4{,}096, which
matches the gradient-step count of BPR-MF. The matched run is the one
cited in the main results table; the short run is reported in
Table~\ref{tab:lgcnmatch} for comparison.

%-----------------------------------------------------------------------------
\section{Evaluation Protocol}
\paragraph{Rating prediction.} MAE $= \frac{1}{|\mathcal{D}|}\sum |r - \hat{r}|$
and RMSE $= \sqrt{\frac{1}{|\mathcal{D}|}\sum (r - \hat{r})^2}$, with
predictions clipped to $[1, 5]$.

\paragraph{Top-N recommendation.} For each user we score all items, mask
out items already in the user's training set (so we never recommend an
item the user has already rated), take the top $K = 10$, and compute
Precision@K, Recall@K, F$_1$@K, and NDCG@K against the user's held-out
test items. Metrics are averaged over all users who have at least one
test item.

\paragraph{Why two models are N/A on MAE/RMSE.} BPR-MF and LightGCN are
trained with the BPR pairwise loss and produce arbitrary real-valued
scores rather than numbers on the 1 to 5 rating scale. Clipping those
scores to $[1,5]$ would produce numerically defined but semantically
meaningless MAE/RMSE values that only reflect the scale mismatch, not
prediction quality. We therefore mark those cells as N/A in the rating
prediction table and evaluate BPR-MF and LightGCN only on the Top-10
task, which is the task those losses were designed for. The six
remaining models are evaluated on both tasks.

%-----------------------------------------------------------------------------
\section{Results}\label{sec:results}

\begin{table}[t]
\centering
\caption{Test-set results on MovieLens-1M. Lower is better for MAE/RMSE; higher is better for P/R/F$_1$/NDCG. BPR-MF and LightGCN do not predict explicit ratings and are only evaluated on Top-$N$.}
\label{tab:results}
\small
\begin{tabular}{lcccccccr}
\toprule
Model & MAE & RMSE & P@10 & R@10 & F$_1$@10 & NDCG@10 & t(s) \\
\midrule
GlobalMean & 0.9331 & 1.1161 & 0.0089 & 0.0022 & 0.0035 & 0.0088 & 0.0000 \\
BiasOnly & 0.7156 & 0.9050 & 0.0525 & 0.0196 & 0.0285 & 0.0492 & 0.4000 \\
ItemKNN & 0.7142 & 0.9308 & 0.0132 & 0.0064 & 0.0086 & 0.0140 & 1.8200 \\
UserKNN & 0.7069 & 0.9182 & 0.0124 & 0.0044 & 0.0065 & 0.0129 & 3.9700 \\
MF & 0.7036 & 0.8873 & 0.1120 & 0.0383 & 0.0571 & 0.1247 & 33.5000 \\
NeuMF & 0.7037 & 0.9043 & 0.0904 & 0.0311 & 0.0463 & 0.0977 & 87.5000 \\
BPR-MF & --- & --- & 0.2861 & 0.1065 & 0.1552 & 0.3202 & 48.9900 \\
LightGCN & --- & --- & 0.1824 & 0.0671 & 0.0981 & 0.2041 & 2270.8200 \\
\bottomrule
\end{tabular}
\end{table}

\subsection{BPR-MF vs.\ MF: Loss Function Matters}
BPR-MF and MF use the \emph{same} scoring function $p_u^\top q_i + b_i$
but different loss functions. BPR-MF beats MF by $2.55\times$ on P@10
($0.286$ vs.\ $0.112$) and $2.57\times$ on NDCG@10 ($0.320$ vs.\ $0.125$).
The reason is that MSE training tries to fit every rating value (including
2s and 3s), which is not directly useful for ranking. BPR training focuses
on the relative order between observed and unobserved items, which is
exactly what the ranking metrics measure.

\subsection{kNN: Good on Rating, Bad on Ranking}
UserKNN and ItemKNN have competitive MAE ($0.707$ and $0.714$ vs.\ MF's
$0.704$), but they perform poorly on ranking. For items a user has not
rated, the kNN prediction tends to fall back toward the user's mean
rating: the neighbours that drive the prediction are mostly items the
user has already rated, and those are masked out in Top-$N$ evaluation.
The remaining scores do not vary enough to produce a good ranking.

\subsection{LightGCN vs.\ BPR-MF}
LightGCN performs worse than BPR-MF on all ranking metrics (P@10:
$0.183$ vs.\ $0.286$; NDCG@10: $0.204$ vs.\ $0.320$). The first run
used 15 epochs at batch size 16{,}384 (about 720 gradient steps),
which is far fewer than BPR-MF's roughly 3{,}900 steps. To rule out the
possibility that LightGCN was simply undertrained, we re-ran with
matched gradient steps (20 epochs, batch size 4{,}096). The numbers do
not change (Table~\ref{tab:lgcnmatch}).

\begin{table}[t]
\centering
\caption{LightGCN under two training budgets on ML-1M. Gradient steps are (epochs $\times$ batches-per-epoch) where a batch covers $|\mathcal{D}_{\text{train}}|$/batch-size samples. The matched configuration uses the same per-batch count and the same number of epochs as our BPR-MF run.}
\label{tab:lgcnmatch}
\small
\begin{tabular}{lrrrrr}
\toprule
Configuration & Epochs & Batch & Grad.\ steps & P@10 & NDCG@10 \\
\midrule
Original (large batch) & 15 & 16\,384 &   $\sim$720  & 0.1826 & 0.2041 \\
Matched (BPR-equivalent) & 20 & 4\,096 & $\sim$3{,}900 & 0.1824 & 0.2041 \\
\midrule
BPR-MF (reference)       & 20 & 4\,096 & $\sim$3{,}900 & \textbf{0.2861} & \textbf{0.3202} \\
\bottomrule
\end{tabular}
\end{table}


The gap is not caused by insufficient training. The likely reasons are:
\begin{enumerate}[leftmargin=*]
\item We used $L = 3$ layers, but $L = 2$ is the optimum reported in the
original paper for ML-1M. We did not sweep this hyperparameter.
\item Although the overall sparsity is 95.16\%, every user has at least 20
ratings, so a plain factor model already has enough direct signal per
user. The extra neighbour-of-neighbour information from graph convolution
may not add much in this setting.
\item Our regularization ($\lambda = 10^{-4}$, Adam defaults) was not tuned.
\end{enumerate}

\subsection{NeuMF}
NeuMF is slightly worse than MF on ranking (NDCG@10: $0.098$ vs.\ $0.125$)
and about the same on rating prediction (MAE: $0.704$ vs.\ $0.704$). The
extra neural layers do not improve results on this dataset.

%-----------------------------------------------------------------------------
\section{Optional Task: Exposure Fairness Audit}\label{sec:fair}
This section addresses the optional trustworthiness task (b), \emph{Fairness
and Unbiases of Recommender Systems}, as listed in the project description.
Beyond ranking accuracy, we check whether each model concentrates its
recommendations on the most popular items. For each model at $K = 10$, we
measure:
\begin{itemize}[leftmargin=*]
\item \textbf{Coverage}: the fraction of items in the catalogue that appear in at least one user's Top-10 list.
\item \textbf{Gini}: the inequality of per-item exposure counts. 0 means every item is exposed equally, 1 means all exposure goes to one item.
\item \textbf{Long-tail Share}: the fraction of recommendation slots that go to items outside the top 20\% most popular items in the training set.
\end{itemize}

\begin{table}[t]
\centering
\caption{Exposure-fairness audit on Top-10 recommendations. Higher Coverage and Long-tail Share are fairer; lower Gini is fairer.}
\label{tab:fairness}
\small
\begin{tabular}{lccc}
\toprule
Model & Coverage@10 $\uparrow$ & Gini $\downarrow$ & Long-tail Share@10 $\uparrow$ \\
\midrule
GlobalMean & 0.1033 & 0.9865 & 0.8727 \\
BiasOnly & 0.0946 & 0.9953 & 0.2597 \\
ItemKNN & 0.8214 & 0.6899 & 0.5568 \\
UserKNN & 0.7573 & 0.7489 & 0.8697 \\
MF & 0.0515 & 0.9937 & 0.0149 \\
NeuMF & 0.4798 & 0.8762 & 0.2573 \\
BPR-MF & 0.2037 & 0.9745 & 0.0047 \\
LightGCN & 0.0217 & 0.9946 & 0.0000 \\
\bottomrule
\end{tabular}
\end{table}

LightGCN places every Top-10 slot on popular items (Long-tail Share
$= 0.000$) and touches only 2.2\% of the catalogue. BPR-MF is similar
(Long-tail Share $= 0.005$, Coverage $= 20.4$\%). The two best-performing
models on ranking accuracy are the two worst on fairness. On the other
hand, UserKNN devotes 87\% of its list to tail items and covers 76\% of
the catalogue, but its ranking accuracy is much lower.

Note that GlobalMean predicts the same score for all items, so its
``top-10'' is effectively a tie-broken arbitrary subset. Its high
Long-tail Share is an artefact of uniform tie-breaking, not a useful
recommendation property, and should not be read as a positive result.

The takeaway is that on ML-1M, the accuracy-optimal models concentrate
recommendations heavily on popular items, and that optimizing only for
NDCG/Precision hides the cost paid by long-tail items. This is a classic
accuracy and fairness trade-off and motivates post-hoc re-ranking or
fairness-aware training objectives, which we leave as future work.

%-----------------------------------------------------------------------------
\section{Web Demo}\label{sec:demo}
The file \texttt{demo/app.py} is a Flask single-page web application
that lets the grader interact with any of the seven trained models.
The user picks a user ID from the MovieLens-1M universe and selects a
model from a dropdown; the server returns the Top-10 recommendation
list alongside that user's highest-rated training items (for
qualitative inspection) and a static table of test-set metrics for
all seven models. Inference is a few milliseconds per user; all
pickled models are loaded once at startup. The demo can be launched
with \texttt{python -m demo.app} and defaults to
\texttt{http://localhost:5000}. Figure~\ref{fig:demo} shows the UI
running on user $1$ with the BPR-MF model selected.

\begin{figure*}[t]
\centering
\includegraphics[width=0.78\textwidth]{../results/demo_screenshot.png}
\caption{Web demo (\texttt{demo/app.py}): test-set metrics for all
seven models (top), live Top-10 recommendation for the selected user
and model (bottom left), and the user's highest-rated training items
for qualitative comparison (bottom right).}
\label{fig:demo}
\end{figure*}

%-----------------------------------------------------------------------------
\section{Reproducibility}
All code is included with this report. The dataset, random seed,
hyperparameters, train/test split, and evaluation protocol are all
fixed and deterministic. The full pipeline is reproduced by:
\begin{verbatim}
bash run_all.sh     # train + eval
python -m demo.app  # launch demo
\end{verbatim}
The script \texttt{run\_all.sh} writes metrics and fairness CSVs to
\texttt{results/}, saves trained checkpoints under \texttt{models/},
and recompiles this report and the slides. The repository README
contains a file-by-file description of every module.

\paragraph{No external recommender libraries.} Per the project policy,
we do not import any packaged recommender-system algorithm. The only
third-party dependencies are numpy, pandas, scipy, scikit-learn,
PyTorch (using only \texttt{nn.Embedding}, \texttt{nn.Linear}, and
\texttt{optim.Adam} as building blocks), tqdm, matplotlib, and flask.
Every similarity formula, every training loop, every evaluation
metric, and every fairness statistic is implemented in this codebase.
No call to Surprise, RecBole, LightFM, implicit, or any similar
library exists anywhere in the project.

%-----------------------------------------------------------------------------
\section{Limitations and Future Work}
Cold-start serving is not implemented: every model is transductive.
Timestamps are unused. Every model is trained with a single seed and
no bootstrap confidence intervals are reported. LightGCN's layer
count is fixed at $L = 3$; sweeping $L \in \{1,2,3,4\}$ is the most
likely fix for the gap to BPR-MF.

\end{document}
